\documentclass[11pt]{article}
\usepackage[utf8]{inputenc}
\usepackage{amsmath,amssymb,amsthm}
\usepackage[margin=1in]{geometry}
\usepackage{hyperref}
\usepackage{xcolor}
\usepackage{enumitem}
\usepackage{newunicodechar}
\newunicodechar{ℝ}{\ensuremath{\mathbb{R}}}
\newunicodechar{⟨}{\ensuremath{\langle}}
\newunicodechar{⟩}{\ensuremath{\rangle}}
\newunicodechar{Γ}{\ensuremath{\Gamma}}
\newunicodechar{·}{\ensuremath{\cdot}}

\title{Tide retrospective:\\mixed quartic-sextic 2D partition (A2)}
\author{Daniel Murfet}
\date{7 May 2026}

\begin{document}
\maketitle

\begin{abstract}
A trivial-mirror tide of A3 (separable-potential abstraction,
Tide~12). For the 2D pure-monomial separable potential
$L(x, y) = x^4/24 + y^6/720$, the partition function and moments
factor as products of the corresponding 1D quartic and sextic
quantities. Closed-form value of the 2D partition function:
$$
  Z_{\mathrm{2D}}(t) = \tfrac{1}{2}(24/t)^{1/4}\Gamma(1/4) \cdot
                       \tfrac{1}{3}(720/t)^{1/6}\Gamma(1/6).
$$
Single session, 153~lines in one new file
(\texttt{Laplace/TwoD/QuarticSextic.lean}), zero \texttt{sorry}, zero
\texttt{axiom}, zero \texttt{native\_decide}. Five public theorems
(partition factorisation, closed form, partition positivity, moment
factorisation, mixed-coordinate covariance vanishing); all are
\texttt{apply}-and-discharge chains over the existing A3 theorems
plus integrability witnesses from the existing pure-monomial
infrastructure. No GPT consult was needed: there is no
candidate-design surface, since the target is the A3 theorem family
instantiated.
\end{abstract}

\section{Setting}

The A-category in the cross-project candidates survey carried three
follow-ups after Tide~12 (the A3 separable-potential abstraction)
landed: A1~(generic even-monomial 1D template), A2~(quartic-sextic
mixed 2D), and A4~(generic mixed even-power 2D moments). A2 is the
smallest and most concrete: a single 2D potential with two
explicitly-named factors, both of which already have closed-form
partition functions in the laplace repo
(\texttt{Laplace.OneD.quartic\_partition} and
\texttt{Laplace.OneD.sextic\_partition}). The candidates survey
estimated $\sim$30~lines.

The seabed at the start of the tide was \texttt{main} at
commit~\texttt{b844d7a}\footnote{Commit prefixes truncated to seven
characters throughout this retrospective; sufficient for uniqueness
across the laplace history.}, with the I2 algebra (Tide
\texttt{tide/gibbscov-algebra}) just merged in. All ingredients for
A2 were on \texttt{main}: the pure-monomial potentials and their
positivity / integrability lemmas in
\texttt{OneD/Quartic.lean} and \texttt{OneD/Sextic.lean}, and the
generic A3 theorem family in \texttt{TwoD/AddSeparable.lean}.

\section{The thing formalised}

Five theorems in the new file
\texttt{Laplace/TwoD/QuarticSextic.lean}, all in the
\texttt{Laplace.TwoD} namespace:

\begin{itemize}[leftmargin=2em]
\item \texttt{partitionFunction\_quarticSextic\_factor}:
$Z_{\mathrm{2D}}(t) = Z_{\mathrm{quartic}}(t) \cdot Z_{\mathrm{sextic}}(t)$
for $t > 0$.

\item \texttt{partitionFunction\_quarticSextic\_eq}: the explicit
closed-form value, obtained by substituting
\texttt{quartic\_partition} and \texttt{sextic\_partition}.

\item \texttt{partitionFunction\_quarticSextic\_pos}: strict
positivity, by \texttt{mul\_pos} on the two pure-monomial positivity
lemmas.

\item \texttt{gibbsExpectation\_quarticSextic\_factor}: for any
separable observable $f(x) \cdot g(y)$ with appropriate
integrability,
$\langle f(x) g(y) \rangle_{\mathrm{2D}, t}
 = \langle f \rangle_{\mathrm{quartic}, t} \cdot
   \langle g \rangle_{\mathrm{sextic}, t}$.

\item \texttt{gibbsCov\_quarticSextic\_fst\_snd\_eq\_zero}: the
covariance of any first-coordinate observable $f(x)$ and any
second-coordinate observable $g(y)$ against the 2D quartic-sextic
Gibbs measure is $0$. The structural statement of independence under
the separable measure.
\end{itemize}

The informal content is exactly what the A3 abstraction says, applied
to one specific pair of potentials. The added value is twofold: (a)
the closed-form value of the 2D partition is now a one-rewrite
substitution rather than a chain through generic infrastructure, and
(b) the moment-factorisation and covariance-vanishing theorems take
\emph{integrability hypotheses against the 1D potentials} directly,
which is what every downstream consumer would need to phrase anyway.

\section{Proof strategy}

Each public theorem is a single \texttt{apply} of the corresponding
A3 generic theorem, with the integrability hypotheses discharged from
the existing pure-monomial infrastructure
(\texttt{quartic\_integrable\_pow\_pot} at $n = 0$ for the quartic
factor; \texttt{sextic\_integrable\_pow\_pot} at $n = 0$ for the
sextic factor). The two integrability discharges are factored into
private helpers
(\texttt{integrable\_exp\_neg\_t\_quartic},
\texttt{integrable\_exp\_neg\_t\_sextic}) so the public proofs are
clean.

For the closed-form theorem, after the factorisation rewrite the
proof is a single \texttt{rw} chain combining the factorisation
lemma with the two pure-monomial closed
forms\footnote{\texttt{partitionFunction\_quarticSextic\_eq};
the chain is \texttt{rw [partitionFunction\_quarticSextic\_factor
ht, quartic\_partition ht, sextic\_partition ht]}.} -- nothing
else needed, because the pure-monomial closed-form expressions are
already exactly the form we want.

\section{Roadblocks and resolutions}

None. The proof rolled through in a single pass with no
re-architecting, no thrashing, no Mathlib-search round. The most
substantive design decision was whether to make the integrability
helpers private to this file or promote them to the AddSeparable
infrastructure file: kept them private here, since they're tied to
the specific (quartic, sextic) instantiation and the
\texttt{integrable\_pow\_pot 0 ht} pattern is generic enough that any
future similar tide can reproduce it in two lines.

\section{What was Mathlib, what was new}

\paragraph{Mathlib lookups.} None new beyond \texttt{ne\_of\_gt} and
\texttt{mul\_pos}.

\paragraph{laplace-side existing material reused.}
\begin{itemize}[leftmargin=2em]
\item The three A3 generic theorems from \texttt{TwoD/AddSeparable.lean}
(Tide~12): the partition factorisation, the moment factorisation
for separable observables, and the mixed-coordinate covariance
vanishing.\footnote{Names: \texttt{partitionFunction\_addSeparable\_factor},
\texttt{gibbsExpectation\_separable\_addSeparable},
\texttt{gibbsCov\_addSeparable\_fst\_snd\_eq\_zero}.}

\item Pure-monomial closed forms and their positivity / integrability
lemmas in \texttt{OneD/Quartic.lean} and
\texttt{OneD/Sextic.lean}.\footnote{Six lemmas:
\texttt{quartic\_partition},
\texttt{quartic\_partition\_pos},
\texttt{quartic\_integrable\_pow\_pot},
\texttt{sextic\_partition},
\texttt{sextic\_partition\_pos},
\texttt{sextic\_integrable\_pow\_pot}.}
\end{itemize}

\paragraph{New material.} Two private integrability helpers and
five public theorems (partition factorisation, partition closed
form, partition positivity, moment factorisation, mixed-coordinate
covariance vanishing). Total: 153~lines, of which $\sim$60 are
docstrings and $\sim$90 are proof body.

\section{Lessons}

\paragraph{Trivial-mirror tides are valuable.}
A2 is the second trivial-mirror tide of the day (after L1 sextic-J-function
on 6~May; today's similar mirrors include G2/G2+ kappa3-affine-anharmonic).
The pattern is: a generic theorem family lands; a strict-improvement
tide instantiates it at specific names. The instantiation is a 50--150
line file, no new mathematics, all proof-engineering value. These
tides are \emph{worth doing} because they (a)~surface the closed-form
values that downstream consumers will want, (b)~front-load the
integrability bookkeeping at the named-instance level rather than
forcing every consumer to repeat it, and (c)~provide a stable named
target for paper-side \texttt{\textbackslash leanref} tags.

\paragraph{No-deliberation tides skip GPT cleanly.}
The skill's "trivial mirror" exemption was used here for the second
time today (after L1 sextic-J-function and G2 anharmonic-affine
mirrors). The criterion is honest: there must be \emph{no candidate
choice} -- the target is forced by the abstraction, the proof is a
single \texttt{apply}-chain, and the integrability hypotheses are
already established. Deliberation overhead would not have surfaced
anything useful. Worth being explicit: the exemption is for tides
where the deliberation surface is genuinely zero, not for tides where
we're \emph{lazy} about deliberation.

\section{Follow-ups}

\begin{itemize}[leftmargin=2em]
\item \emph{A4 -- mixed even-power 2D moments.} The natural sequel:
the closed form of $\langle x^{2j} y^{2k} \rangle_{\mathrm{2D}, t}$
under the same potential, by applying
\texttt{gibbsExpectation\_quarticSextic\_factor} to
$f = x^{2j}$, $g = y^{2k}$, then substituting
\texttt{quartic\_expected\_value\_even} and
\texttt{sextic\_expected\_value\_even}. Estimated $\sim$30~lines.

\item \emph{A1 -- generic even-monomial 1D template
$L = x^{2k}/(2k)!$.} The natural \emph{abstraction} step: unify the
quartic, sextic, and (now) generic-$k$ J-function machinery under a
single parametric family. Estimated $\sim$250--350~lines per the
candidates survey.

\item \emph{Other 2D mixed pure-monomial pairs.} Quartic-quartic,
sextic-sextic, $x^{2j}/( (2j)!)$ paired with $y^{2k}/((2k)!)$ for
various $(j, k)$. Each is a 30-line corollary of A3 once the named
1D pure-monomial families exist. Best deferred until A1 lands and
gives the parametric 1D pattern.

\item \emph{$\backslash$leanref tag for the primer.} The
2D mixed quartic-sextic example may appear in the susceptibility
primer's exposition; if so, tag with
\texttt{\textbackslash laplaceCommit} pinned to the merge commit.
Process tide.
\end{itemize}

\section*{Acknowledgements}

This Tide consumed only material on \texttt{main} at the time of
branching (commit~\texttt{b844d7a}, post-I2-merge). It is genuinely
non-rivalrous with the half-dozen other tides running concurrently:
A2 only writes a new file, doesn't touch any existing module, and
its only dependency on today's earlier work is the I2 algebra
(transitively, via \texttt{Laplace.gibbsExpectation\_const} inside
the A3 theorems' proofs).

GPT-5.5~Pro was not consulted: the trivial-mirror exemption applies,
since there are no candidate-design choices and the proof is a
direct \texttt{apply}-chain over the A3 theorem family. This is the
second use of that exemption today (the first was the L1 sextic-J-function
mirror tide on 6~May).

\end{document}
